% Part: lambda-calculus
% Chapter: introduction
% Section: lambda-definability

\documentclass[../../../include/open-logic-section]{subfiles}

\begin{document}

\olfileid[fr]{lam}{int}{rep}
\olsection{Fonctions arithmétiques \texorpdfstring{$\lambd$}{λ}-définissables}

Comment le lambda-calcul peut-il servir de modèle
de calcul? À première vue, le sens même de cette
question n'est pas clair. Pour parler de calculabilité
sur les nombres naturels, il nous faut une représentation
appropriée de ces nombres. En voici une qui fonctionne
remarquablement bien.

\begin{defn}
Pour chaque entier naturel~$n$, le \emph{numéral de Church}
$\num{n}$ est le lambda-terme
$\lambd[x][\lambd[y][(x(x(x(\dots x(y)))))]]$,
où le corps contient $n$ occurrences de~$x$. Pour $n=0$, il est
simplement~$y$.
\end{defn}

Les termes $\num{n}$ sont des « itérateurs » :
appliqué à~$f$, $\num{n}$ renvoie la fonction qui
associe $f^n(y)$ à~$y$. Chaque numéral est sous
forme normale. Nous pouvons maintenant préciser
ce que signifie, pour un lambda-terme, « calculer »
une fonction sur les nombres naturels.

\begin{defn}
Soit $f(x_0, \dots, x_{k-1})$ une fonction partielle
à $k$ arguments de $\Nat^k$ dans~$\Nat$. Nous disons
qu'un lambda-terme~$F$ \emph{définit}~$f$ si, pour
toute suite d'entiers naturels $n_0$, \dots,~$n_{k-1}$,
\[
F\, \num{n_0}\, \num{n_1} \dots \num{n_{k-1}}
\red \num{f(n_0, n_1, \dots, n_{k-1})}
\]
lorsque $f(n_0, \dots, n_{k-1})$ est défini;
si cette valeur n'est pas définie, le terme
$F\, \num{n_0}\, \num{n_1} \dots \num{n_{k-1}}$
n'a pas de forme normale.
\end{defn}

\begin{thm}
\ollabel{thm:lambda-def}
Une fonction~$f$ est partielle calculable si
et seulement si elle est définie par un lambda-terme.
\end{thm}

\begin{explain}
Ce théorème est quelque peu frappant. Comme modèle
de calcul, le lambda-calcul est fort simple : ses
seules opérations sont l'abstraction lambda et
l'application!{} Pourtant, avec des moyens aussi
réduits, il permet de réaliser toute procédure de
calcul.
\end{explain}

\begin{explain}
Le terme appliqué reçoit exactement $k$ numéraux. Pour une
entrée où $f$ est définie, il faut pouvoir atteindre le
numéral de la valeur. Pour une entrée où $f$ est indéfinie,
aucune suite de réductions ne doit aboutir à une forme normale.
\end{explain}

\end{document}
